\section{Introduction}\label{sec:intro}

A head-mounted display shows a world that must respond to the head that wears it. When the user turns,
the new view must appear within a few tens of milliseconds, or the image lags the body and comfort
fails \cite{W2784711283}. On a standalone device this deadline is met locally. Once the scene, the
computation or the other participants are elsewhere, the network becomes part of the perceptual loop.
Extended reality (XR), the umbrella term for virtual, augmented and mixed reality, therefore asks more of
a network than video or the web. A 360-degree or volumetric stream needs hundreds of megabits per second
if it is sent whole \cite{W3038754372,W3096564590}. Remote rendering must finish a frame within the
motion-to-photon (MTP) budget \cite{W2784711283}. The device sends pose, camera frames and captured media
upstream all the time. When several people share a world, the network must also keep them consistent
with each other \cite{W2561715836}.

Ten years of research have produced many answers to these demands, and they are scattered across
communities that seldom read each other. Multimedia-systems researchers built viewport-adaptive
streaming. Mobile-computing researchers offloaded AR recognition to the edge \cite{W3133913697}. Wireless
researchers framed XR as a joint allocation of radio, computing and energy. Haptics researchers pursued
the tactile internet \cite{W2811040154}. Visionaries placed all of it inside 6G and the metaverse
\cite{W4312940001,W4390204075}. Existing surveys follow these boundaries. Each covers one medium, such as
360-degree video \cite{W3038754372,W2971680322,W2917103012}, point clouds \cite{W3096564590} or haptics
\cite{W3088582486}, or one technology, such as edge computing \cite{W3133913697,W2902319104}, ultra-low
latency \cite{W4362585626} or the 6G metaverse \cite{W4312940001,W4315497913}. None measures the field as
a whole, and none asks how well its claims are tested.

This review does both. It condenses the \emph{Networked XR Atlas}, a corpus of \nxrIncluded{} papers
published between October 2016 and October 2026. The atlas was screened from \nxrHarvested{} records
with a validated procedure, labelled on five dimensions and read in depth where it mattered most
(Section~\ref{sec:method}). The atlas supports three kinds of statement: counts and trends over the whole
corpus, syntheses of the papers that shape each topic, and gaps that have been checked against a text
search of every included paper.

\paragraph{Organising idea.} We do not organise the field by technology or by medium. We organise it by
\emph{function}: the six things any networked XR system must do, whatever its medium and access network
(Section~\ref{sec:framework}). A system must \emph{encode} the scene into something that can be sent
(Section~\ref{sec:encode}). It must \emph{anticipate} what the user will look at or do next, so that the
right content arrives in time (Section~\ref{sec:anticipate}). It must decide where to \emph{compute}, on
the device, at the edge or in the cloud, and how to hide the delay this adds (Section~\ref{sec:compute}).
It must \emph{provision} the network with radio, scheduling, transport and architecture that supply the
rate, delay and reliability required (Section~\ref{sec:provision}). It must \emph{connect} users with each
other and with remote machines, through shared state, touch and telepresence (Section~\ref{sec:connect}).
And it must \emph{assure} that the result is good enough and safe for the people who use it
(Section~\ref{sec:assure}). Each function poses its own design question, and the answers can be compared
on the same terms.

\paragraph{Contributions.} The review makes four contributions.
\begin{enumerate}[leftmargin=*,label=(\arabic*)]
  \item A \textbf{functional taxonomy} of networked XR that places the \nxrNCatProb{} networking problems
        and \nxrNCatCont{} content types of the atlas codebook into six functions, each defined by a design
        question (Section~\ref{sec:framework}, \cref{fig:architecture,tab:framework}).
  \item A \textbf{quantitative account of how the field has changed}: its growth, the move from
        delivering content to allocating resources, the replacement of 360-degree video by generic XR
        traffic as the reference workload, and the decline of evaluation on real systems
        (Section~\ref{sec:landscape}).
  \item A \textbf{synthesis of each function} that separates what is established, what is contested and
        what is assumed, with the measured results that support each claim
        (Sections~\ref{sec:encode}--\ref{sec:assure}).
  \item \textbf{Nine cross-cutting gaps} and a research agenda ordered by dependency: which projects
        produce the data and models that the others need (Sections~\ref{sec:synthesis}
        and~\ref{sec:agenda}).
\end{enumerate}
